% ============================================================================
%  Chapter 5 --- Data Strategy
% ============================================================================
\chapter{Data Strategy}
\label{ch:data}

\section{Introduction}

The credibility of any AI system rests on the quality and authenticity of the data it is trained on. This chapter presents the project's data strategy: the real data sources from Tunisie Telecom and Huawei, the ingestion pipeline, the anonymisation process, the schema mapping, and the dual-mode architecture that allows the platform to operate on both real and synthetic data.

\section{Data Sources}

The platform ingests data from two primary sources and maintains a synthetic fallback:

\begin{table}[H]
\centering
\caption{Data sources and their characteristics.}
\label{tab:data-sources}
\begin{tabularx}{\textwidth}{@{}llllX@{}}
\toprule
\textbf{Source} & \textbf{Domain} & \textbf{Format} & \textbf{Status} & \textbf{Content} \\
\midrule
Tunisie Telecom & BSS & CSV / Excel & Real (provided via Huawei) & Subscriber profiles, ARPU, DOU, churn, plan types, revenue, usage \\
Tunisie Telecom & OSS & CSV / JSON & Real (provided via Huawei) & Cell-level KPIs: throughput, latency, packet loss, active users, RSRP \\
Huawei SmartCare & CEM & CSV / JSON & Real (export) & KQI, CEI scores, demarcation results, experience alerts \\
Synthetic Gen. & OSS + BSS & In-memory & Fallback / demo & Generated by pipeline-worker when no real data is available \\
\bottomrule
\end{tabularx}
\end{table}

\subsection{Why Real Data Matters}

The shift from synthetic-only to real operator data fundamentally changes the project's value:

\begin{itemize}
  \item \textbf{Before (synthetic):} Models were trained on artificially generated data with injected patterns. The anomalies were designed, the correlations were scripted, and the evaluation was circular --- the model could only ``discover'' what was explicitly injected.
  
  \item \textbf{After (real):} Models are trained on genuine network behaviour from a live Tunisian operator. The anomaly patterns are organic (not injected), the correlations are discovered (not scripted), and the evaluation metrics (precision, recall, F1) measure real detection capability.
  
  \item \textbf{Defence-proof:} When the jury asks "how do you know your model works?", the answer is: "it was trained on anonymised production data from Tunisie Telecom and validated against known network events."
\end{itemize}

\section{OSS Data Schema}

The following schema defines the expected structure of real OSS data from Tunisie Telecom. The exact number of records will depend on the dataset provided.

\begin{table}[H]
\centering
\caption{Expected OSS data schema (real TT data).}
\label{tab:oss-schema}
\begin{tabularx}{\textwidth}{@{}llcX@{}}
\toprule
\textbf{Source Column} & \textbf{Internal Field} & \textbf{Type} & \textbf{Description} \\
\midrule
cell\_id / eNodeB\_id     & \texttt{cell\_id}           & string  & Cell tower identifier \\
throughput\_dl\_mbps       & \texttt{throughput\_mbps}    & float   & Downlink throughput (Mbps) \\
latency\_rtt\_ms           & \texttt{latency\_ms}         & float   & Round-trip latency (ms) \\
packet\_loss\_pct          & \texttt{packet\_loss\_pct}   & float   & Packet loss percentage \\
active\_ue\_count          & \texttt{active\_users}       & int     & Number of active UEs on cell \\
rsrp\_dbm                  & \texttt{signal\_rsrp\_dbm}   & float   & Reference Signal Received Power (dBm) \\
region / site\_name        & \texttt{region}              & string  & Gouvernorat or site name \\
timestamp                  & \texttt{ts}                  & datetime & Measurement timestamp \\
\bottomrule
\end{tabularx}
\end{table}

\section{BSS Data Schema}

The BSS schema captures subscriber-level business metrics. The inclusion of APPU and DOU as explicit fields was directed by the supervisor to align with Huawei's BSS analytics framework.

\begin{table}[H]
\centering
\caption{Expected BSS data schema (real TT data).}
\label{tab:bss-schema}
\begin{tabularx}{\textwidth}{@{}llcX@{}}
\toprule
\textbf{Source Column} & \textbf{Internal Field} & \textbf{Type} & \textbf{Description} \\
\midrule
subscriber\_id / MSISDN & \texttt{subscriber\_id}  & string  & Anonymised subscriber identifier \\
operator                 & \texttt{operator}         & string  & Operator name (TT, Ooredoo, Orange) \\
line\_type               & \texttt{line\_type}       & string  & \texttt{prepaid} or \texttt{postpaid} \\
plan / forfait code      & \texttt{plan}             & string  & Subscription plan identifier \\
revenue (TND)            & \texttt{revenue\_tnd}     & float   & Monthly revenue in Tunisian Dinars \\
data\_usage\_gb           & \texttt{data\_used\_gb}   & float   & Monthly data consumption (GB) \\
voice\_minutes            & \texttt{voice\_min}       & float   & Monthly voice usage (minutes) \\
sms\_count                & \texttt{sms\_count}       & int     & Monthly SMS count \\
churn\_indicator          & \texttt{churn\_risk}      & float   & Churn risk score $\in [0, 1]$ \\
\textbf{appu\_tnd}       & \texttt{appu\_tnd}        & float   & Average Purchase Per User (TND) \\
\textbf{dou\_gb}         & \texttt{dou\_gb}          & float   & Data of Use --- monthly consumption (GB) \\
region / gouvernorat     & \texttt{region}            & string  & Subscriber's region \\
serving\_cell             & \texttt{serving\_cell}    & string  & Current serving cell (for OSS join) \\
timestamp / period       & \texttt{ts}               & datetime & Observation period \\
\bottomrule
\end{tabularx}
\end{table}

\section{Data Ingestion Pipeline}

% ── Figure: Data Ingestion Flow ───────────────────────────────────────────────
\begin{figure}[H]
\centering
\begin{tikzpicture}[node distance=1.2cm and 2cm]

  % Source files
  \node[external, minimum width=2.5cm] (csv) {CSV / Excel\\from TT};
  \node[external, minimum width=2.5cm, below=0.8cm of csv] (json) {JSON / API\\from Huawei};
  
  % Ingestion service
  \node[service=cemteal, right=2cm of csv, yshift=-0.5cm, minimum width=3.5cm, minimum height=2.2cm] (ingest) {
    \textbf{Data Ingest}\\[2pt]
    {\scriptsize Column mapping}\\
    {\scriptsize Type validation}\\
    {\scriptsize Anonymisation}\\
    {\scriptsize Format normalisation}
  };
  
  % Schema validation
  \node[service=ossblue, right=2cm of ingest, minimum width=3cm] (validate) {
    \textbf{Schema}\\
    \textbf{Validation}\\
    {\scriptsize Required fields}\\
    {\scriptsize Type checks}
  };
  
  % Raw layer
  \node[datastore=datalake, right=2cm of validate] (rawlake) {
    \textbf{Raw Layer}\\{\scriptsize MinIO / OBS}
  };
  
  % Arrows
  \draw[myarrow=cemteal!70] (csv) -| (ingest);
  \draw[myarrow=cemteal!70] (json) -| (ingest);
  \draw[myarrow=cemteal!70] (ingest) -- (validate)
    node[arrlabel, midway, above] {Mapped data};
  \draw[myarrow=ossblue!70] (validate) -- (rawlake)
    node[arrlabel, midway, above] {Validated};
    
  % Synthetic fallback
  \node[external, below=1.5cm of ingest, minimum width=3.5cm] (synth) {
    Synthetic Generator\\{\scriptsize (fallback if no real data)}
  };
  \draw[myarrow=gray, dashed] (synth) -- (validate)
    node[arrlabel, midway, below right] {\scriptsize fallback};
    
\end{tikzpicture}
\caption{Data ingestion pipeline: real TT data is mapped, anonymised, validated, and stored in the raw lake layer. Synthetic generation is the fallback.}
\label{fig:data-ingestion}
\end{figure}

\subsection{Column Mapping}

The ingestion service performs column mapping from TT's native schema to the platform's internal schema. This mapping is configurable --- when TT provides data with different column names (e.g., \texttt{dl\_throughput} instead of \texttt{throughput\_dl\_mbps}), the mapping configuration is updated without code changes.

\subsection{Dual-Mode Operation}

The pipeline-worker operates in two modes, controlled by the environment variable \texttt{DATA\_SOURCE}:

\begin{itemize}
  \item \texttt{DATA\_SOURCE=real} (default): The worker reads real TT data files from a mounted volume (\texttt{/data/tt-import/}), applies anonymisation and column mapping, then proceeds with the standard 5-phase pipeline.
  \item \texttt{DATA\_SOURCE=synthetic}: The worker generates synthetic data using the built-in generators (retained from the earlier project phases). This mode is used for demos, testing, and when real data is not available.
\end{itemize}

\section{Anonymisation}

All personally identifiable information (PII) is removed or pseudonymised before the data enters the platform:

\begin{table}[H]
\centering
\caption{Anonymisation steps applied to incoming TT data.}
\label{tab:anonymisation}
\begin{tabularx}{\textwidth}{@{}llX@{}}
\toprule
\textbf{Field} & \textbf{Method} & \textbf{Rationale} \\
\midrule
MSISDN / subscriber\_id & SHA-256 hash with salt & Enables record linkage without exposing phone numbers \\
Name / address / NIN     & Strip entirely & Not required for analysis \\
Geographic precision     & Gouvernorat-level only & No GPS coordinates or precise addresses \\
Cell IDs                 & Optional opaque mapping & Map to anonymised IDs if TT requires \\
Temporal precision       & Minute-level preserved & Required for meaningful OSS--BSS correlation \\
\bottomrule
\end{tabularx}
\end{table}

The anonymisation is applied \emph{during ingestion}, before data reaches the raw layer. The raw layer thus contains already-anonymised data --- no PII exists anywhere in the platform.

\section{Feature Engineering}

Once raw data is ingested and stored, the pipeline performs feature engineering to transform it into ML-ready features.

\subsection{OSS Feature Engineering}

Each raw OSS record is enriched with computed fields:

\begin{itemize}
  \item \textbf{latency\_severity:} Categorical severity based on latency thresholds (\texttt{normal}, \texttt{warning}, \texttt{critical}).
  \item \textbf{throughput\_category:} Categorical label based on throughput ranges (\texttt{poor}, \texttt{fair}, \texttt{good}, \texttt{excellent}).
  \item \textbf{load\_factor:} $\text{active\_users} / \text{cell\_capacity}$ --- normalised cell utilisation.
  \item \textbf{qos\_score:} Composite quality score combining latency, throughput, and packet loss into a single $[0, 1]$ indicator.
\end{itemize}

For the SLA risk model, records are aggregated into \textbf{time windows} (per region, per time period) to produce 9 aggregate KPI features:

\begin{enumerate}
  \item \texttt{mean\_throughput\_mbps}
  \item \texttt{std\_throughput\_mbps}
  \item \texttt{mean\_latency\_ms}
  \item \texttt{std\_latency\_ms}
  \item \texttt{max\_latency\_ms}
  \item \texttt{mean\_packet\_loss\_pct}
  \item \texttt{max\_packet\_loss\_pct}
  \item \texttt{mean\_active\_users}
  \item \texttt{mean\_signal\_rsrp\_dbm}
\end{enumerate}

\subsection{BSS Feature Engineering}

Each BSS record is enriched with:

\begin{itemize}
  \item \textbf{arpu\_category:} \texttt{low} ($<$ 10 TND), \texttt{mid} ($<$ 40 TND), \texttt{high} ($\geq$ 40 TND)
  \item \textbf{data\_intensity:} Consumption brackets based on \texttt{dou\_gb}
  \item \textbf{churn\_bucket:} \texttt{safe}, \texttt{watch}, \texttt{risk}, \texttt{critical} based on \texttt{churn\_risk}
\end{itemize}

The BSS feature space for the revenue anomaly model includes 7 features:
\begin{enumerate}
  \item \texttt{revenue\_tnd}
  \item \texttt{data\_used\_gb}
  \item \texttt{voice\_min}
  \item \texttt{sms\_count}
  \item \texttt{churn\_risk}
  \item \texttt{appu\_tnd} (new --- Average Purchase Per User)
  \item \texttt{dou\_gb} (new --- Data of Use)
\end{enumerate}

\section{Data Quality and Validation}

Before processing, the ingestion service applies validation rules:

\begin{itemize}
  \item \textbf{Required fields:} All mandatory columns must be present (nulls are acceptable for optional fields).
  \item \textbf{Type checks:} Numeric fields must parse as floats/ints; timestamps must parse as ISO 8601 or common date formats.
  \item \textbf{Range checks:} Throughput $> 0$, latency $> 0$, packet loss $\in [0, 100]$, RSRP $\in [-140, -40]$ dBm.
  \item \textbf{Deduplication:} Exact duplicate records (same cell\_id + timestamp + all metrics) are removed.
\end{itemize}

Records that fail validation are logged but not silently dropped --- the pipeline reports the number of rejected records per run.

\section{Data Lineage and Traceability}

Every dataset produced by the pipeline is registered in the \texttt{dataset\_registry} table with:
\begin{itemize}
  \item The \texttt{run\_id} that produced it
  \item The \texttt{layer} (raw, processed, curated)
  \item The \texttt{object\_key} pointing to the MinIO/OBS object
  \item The \texttt{row\_count}
  \item The \texttt{schema\_version}
\end{itemize}

This enables full traceability: any anomaly or SLA risk score can be traced back to the exact raw data and pipeline run that produced it.

\section{Chapter Summary}

This chapter detailed the data strategy: real data from Tunisie Telecom and Huawei flows through an ingestion pipeline that anonymises, maps, validates, and stores data in a three-layer lake. The platform operates in dual mode (real or synthetic), ensuring robustness across deployment scenarios. The next chapter describes the AI models trained on this data.
